% Глава 1. Материал: docs/requirements/audience.md, docs/research/literature-review.md,
% docs/research/preregistration.md, reports/ (SRC-038).
\chapter{Введение}\label{ch:intro}

\section{Контекст}

Чтение Корана вслух --- ежедневная практика многих мусульман, и правильное произношение --- часть самой
практики. Для арабоязычных главная задача --- выучить правила чтения; для всех остальных путь начинается
гораздо раньше, с букв и звуков за ними. В России и Средней Азии у этого пути долгая и особая история.
В конце XIX века реформаторы-джадиды заменили старый метод, при котором ученики годами заучивали
названия букв, звуковым: буква изучалась как звук, а из звуков небольшими шагами строились слоги
\cite{SRC-016,SRC-017}. Букварь Ахмад-Хади Максуди «Муаллим ас-Сани» («Второй учитель») перенёс этот
метод на чтение Корана \cite{SRC-015}. Больше века спустя его по-прежнему издаёт Духовное управление
мусульман Татарстана и официально распространяют в Казахстане \cite{SRC-040,SRC-020}.

По букварю никогда не учились без учителя. Урок следует циклу \emph{таляккы}: учитель читает, ученик
повторяет, учитель исправляет каждую оговорку, и ученик идёт дальше только после безошибочного чтения
\cite{SRC-021,SRC-058}. Цикл действенный, но он привязывает каждого ученика ко времени учителя. Взрослые,
которые занимаются после работы, живут далеко от медресе или стесняются читать при других, упражняются
между уроками одни, и никто не говорит им, что \ar{ح} превратилась в русское [х], а долгий гласный стал
кратким.

Два обстоятельства позволяют спросить, не может ли машина закрыть часть этого разрыва.
Распознавание речи для коранического арабского дошло до уровня фонем: открытая модель Quran Muaalem
распознаёт фонемы вместе с десятью артикуляционными признаками и на чтении экспертов ошибается примерно
в половине процента фонем \cite{SRC-002}. А исследования компьютерного обучения произношению показали,
что автоматическая обратная связь даёт средний положительный эффект, наибольший тогда, когда она
нацелена на отдельные звуки и прямо называет ошибку \cite{SRC-025}, --- ровно то, что делает букварь.

\section{Проблема}

Трудность в том, кто говорит. Muaalem обучена только на чтении мужчин-экспертов. На записях с настоящими
ошибками доля ошибок по фонемам у неё в десять раз выше, чем на собственном тесте, а некоторые правила
распознаются с F1 ниже половины \cite{SRC-002}. Открытые бенчмарки чтения Корана содержат беглое чтение
аятов и не содержат изучающих арабский как второй язык \cite{SRC-004,SRC-031,SRC-043}. Систематические
обзоры предупреждают, что распознаватели дают ложные замечания на речи с акцентом, и советуют оставлять
в контуре учителя или партнёра \cite{SRC-026}. Русскоязычный взрослый, впервые медленно читающий
отдельный слог, далёк от обучающих данных модели, и никто не знает, как часто модель будет ошибаться
насчёт такого ученика, как часто отклонит верное чтение и помогут ли её объяснения.

Вторая проблема касается того, как такая система создаётся. Программное обеспечение этой работы написано
одним автором вместе с ИИ-агентом. Агенты продуктивны, но известно, что они добиваются прохождения
собственных тестов, меняя сами тесты \cite{SRC-027}, а письменные инструкции теряют силу за долгую сессию
\cite{SRC-076}. Исследовательскому ПО, чьи числа попадают в диссертацию, нужны гарантии сильнее
файла с инструкциями.

\section{Цель и задачи}

Цель работы --- \emph{спроектировать тренажёр чтения, точно следующий} «Муаллим ас-Сани», \emph{и измерить,
насколько открытая модель распознавания фонем может взять на себя роль слушающего учителя для
русскоязычных начинающих, не преувеличивая её возможностей.}

Задачи:
\begin{enumerate}[label=З\arabic*.]
  \item перевести букварь в цифровую программу в порядке самой книги так, чтобы каждое упражнение
    прослеживалось до страницы и строки конкретного издания;
  \item создать тренажёр, который оценивает чтение на уровне фонем, объясняет ошибки по-русски для
    ученика «с нуля» и сохраняет полный след решения по каждой попытке;
  \item построить инфраструктуру исследования: согласие и псевдонимы, разметку по фонемам двумя
    преподавателями, автоматический расчёт метрик и предрегистрированный план анализа;
  \item провести закрытый пилот с 20--40 учениками в течение 4--6 недель, включая внутрисубъектный
    эксперимент с обратной связью и отложенный тест;
  \item ответить на исследовательские вопросы с размерами эффекта и доверительными интервалами;
  \item описать обвязку ИИ-агента, в которой создана система, и оценить её как кейс-стади.
\end{enumerate}

\section{Исследовательские вопросы}\label{sec:rq}

\begin{rqbox}
\textbf{RQ-1.} Насколько автоматическая проверка --- Muaalem, выравнивание по правилам и пороги оценки ---
согласуется с квалифицированными преподавателями на чтении начинающих по «Муаллим ас-Сани»?\\[4pt]
\textbf{RQ-2.} Какие ошибки чтения делают русскоязычные начинающие и какие из ошибок, найденных
тренажёром, принадлежат ученику, а не модели?\\[4pt]
\textbf{RQ-3.} Улучшают ли явные объяснения на русском следующую попытку и отложенное запоминание по
сравнению с простым сигналом «верно / неверно»?\\[4pt]
\textbf{RQ-4.} Какие известные сбои разработки с ИИ-агентами предотвращает или делает видимыми механически
обеспеченная обвязка и сколько вмешательств человека она требует?
\end{rqbox}

RQ-1--RQ-3 касаются тренажёра, RQ-4 --- того, как он создан. Семь подтверждающих гипотез к RQ-1--RQ-3
перечислены в разделе~\ref{sec:method-hyp}; всё остальное подаётся как разведочный анализ.

\section{Подход}

Работа выполнена в парадигме design science research \cite{SRC-100,SRC-101}. В ней создаются два артефакта,
и каждое утверждение оценивается подходящим методом. Тренажёр (артефакт A) оценивается измерением
относительно преподавателей, описательным анализом ошибок и пилотным экспериментом. Обвязка разработки
(артефакт B) оценивается как встроенное кейс-стади \cite{SRC-102}. Требования записаны как сценарии
атрибутов качества \cite{SRC-047}, метрики выведены методом «цель --- вопрос --- метрика» \cite{SRC-048},
а план анализа предрегистрирован и зафиксирован отпечатком до сбора данных.

\section{Вклад}

\begin{enumerate}
  \item \textbf{Эмпирический.} Первое измерение открытой модели распознавания фонем Корана на букварном
    чтении начинающих, для которых арабский не родной, относительно оценки преподавателей, включая то,
    как часто она ошибочно отклоняет верное чтение.
  \item \textbf{Данные.} Небольшой набор чтения начинающих, собранный с согласия, с разметкой по фонемам
    от двух преподавателей, подготовленный к публикации в формате Iqra'Eval, если позволят этика и
    согласия.
  \item \textbf{Проектный.} Тренажёр, функции которого выведены из дидактических принципов Максуди и
    исследований обратной связи, с открытой архитектурой: модель, базу данных и мессенджер можно
    заменить, не трогая предметную логику.
  \item \textbf{Методический.} Доказательное описание обвязки ИИ-агента для исследовательского ПО: что
    она предотвращает, чего стоит и где даёт сбой.
\end{enumerate}

\section{Границы и ограничения}

Программа охватывает всю книгу (с.~10--85), но оценка сосредоточена на уроках букв и огласовок
(с.~10--24): с них начинают и в них больше всего ошибок. Правила таджвида вроде \emph{ихфа} и
\emph{идгам}, другие родные языки и дообучение модели оставлены для дальнейшей работы. Тренажёр
задуман как дополнение к учителю; работа не утверждает, что он может учителя заменить, --- и сама книга
говорит, что не может.

\section{Структура работы}

Глава~\ref{ch:review} рассматривает букварь и его метод, учеников, исследования обратной связи по
произношению, автоматическую оценку чтения Корана и программную инженерию с ИИ-агентами.
Глава~\ref{ch:method} излагает методологию: требования, метрики, четыре исследования, предрегистрацию и
угрозы валидности. Глава~\ref{ch:impl} описывает реализацию --- от программы по книге до конвейера оценки,
исследовательских данных и обвязки разработки. Глава~\ref{ch:eval} представляет и обсуждает оценку, а
глава~\ref{ch:concl} завершает работу ответами и направлениями для докторантуры.
